% Computer Graphics and VR chapter (G). Compile from report/. Tables: report/tables/e10.tex and
% e11.tex, written by `python -m experiments.report_tables` from bench/results (rule 11).
% Techniques are described from the module headers in viewer/src and viewer/shaders; every
% screenshot must come from the real viewer, captioned with commit and data source.
% Bibliography keys: mcguire2013wboit, rong2006jfa, correll2018vsup (CLAUDE.md §16.1 lists VSUP;
% add the WBOIT and JFA papers to §16.1 before citing them, or cite them as the module headers
% do, by author and year, with the reviewer's approval).

\chapter{Seeing a certified field}
\label{ch:cgvr}

The viewer shows a 100~mm four-layer board at 1~m width on a desktop or in a WebXR headset. It
renders the temperature field the server sends, lets the user edit the design, and makes the
reliability of what is shown visible: a prediction, a field being verified, and a field verified
within a guaranteed band are drawn differently and labelled. Three.js provides the scene graph and
the WebXR session; every shader and technique below was written for this project
(\texttt{viewer/src}, \texttt{viewer/shaders}).

\section{Architecture}
The server and the viewer talk through versioned JSON contracts over a WebSocket
(\texttt{contracts/schemas}): \texttt{query} and \texttt{design} from the viewer, \texttt{pred},
\texttt{iterate} and \texttt{verified} from the server, each field a base64 float16 array
$[\text{layer},y,x]$. Every message is validated against the same schemas in Python and in
JavaScript before anything is drawn (invariant B6). Parsing, validation and decoding of a
$256\times256\times4$ field take tens of milliseconds, so they run in a Web Worker; the main
thread only copies the decoded half-float bits into one R16F texture per layer.

\section{Rendering techniques}
\paragraph{Layer peel and through-thickness section (T1.1).} The four layer planes separate on a
slider, or in VR by the change in distance between the two controllers while both grips are held,
with exponential smoothing so tracking jitter never makes the stack jump. Thermal vias are drawn as
instanced cylinders, one per via per gap, whose glow follows the vertical heat flux
$\nu k_{\mathrm{Cu}} A_v(\theta_l-\theta_m)/d_{lm}$ they carry, tone-mapped by a Reinhard curve.
The FR-4 dielectric is a translucent slab composited with weighted blended order-independent
transparency (McGuire and Bavoil, 2013): an additive accumulation pass into an RGBA16F target, a
multiplicative revealage pass into an R8 target and a full-screen composite, so overlapping
translucent layers need no sorting. A vertical cut plane shows the section through the stack by
interpolating linearly between the layer textures.

\paragraph{Colour and isolines (T1.3).} Temperature maps through a 256-entry lookup table
(inferno by default, ironbow and white-hot) shared by all layers and the colour bar. Isolines are
drawn analytically in the fragment shader: the distance to the nearest level, divided by
\texttt{fwidth} of the level coordinate, is a distance in screen pixels, so lines keep a constant
width at any distance and fade out where levels crowd together instead of aliasing. A probe
raycasts the planes and reports $25+\theta$~\textdegree C at the cell under the pointer.

\paragraph{A realistic board (T1.4).} A jump-flooding pass (Rong and Tan, 2006) on the top copper
mask gives a signed distance field in $\log_2 N$ full-screen passes; a smooth bevel profile of that
distance yields a normal map, so copper pours read through the green solder mask with a raised
edge. Copper, mask and silkscreen are physically based materials lit by image-based lighting,
and every chip of the design record gets an instanced package exactly on the footprint where the
solver deposits its heat, so the picture never shows a heat source the physics does not have.

\section{Trust and anytime verified rendering (T1.6, T1.7)}
The defining feature of the viewer is that it shows how far the picture can be trusted, using only
what the server certifies. A status badge is driven exclusively by the type of the last valid
message (\texttt{viewer/src/trust.js}):
\begin{itemize}
\item \textbf{Predicted} (amber): a network prediction. With the live gate, the badge says whether
the gate accepted it ($B = q_t s \le \tau$; the prediction is used and nothing is solved), flagged
it, or selected it for a random audit. While a chip is dragged, previews are never verified.
\item \textbf{Verifying} (blue): the server streams certified iterates of the exact solver,
warm-started from the prediction (Chapter~\ref{ch:pagp}). Each iterate replaces the drawn field,
so the board visibly converges, and carries a guaranteed band: a number $c_k$ such that the true
temperature of \emph{every} cell lies within $\pm c_k$ of the value drawn. The band includes the
float16 rounding of the transmitted field, so it is rigorous for what the user actually sees.
\item \textbf{Verified} (green): the band is below the target (0.1~K), with the iterations and
time taken. A board without a certificate is shown as solved with no band claimed.
\end{itemize}
A live design panel turns the viewer into the design tool the project is about: the user picks a
board, moves a chip in millimetre steps (each step a preview) and commits; the server's gate
decides, the audit coin is drawn after the commit, and flagged or audited designs stream their
verification into the same view.

\section{Results}
\paragraph{Frame rate (E10).} Table~\ref{tab:e10} records 30 seconds of frame times on the
desktop while a verification stream plays. The viewer holds the display's refresh rate, meeting
the 60~fps desktop target; the frame rate in a headset was not measured.

% generated by experiments/report_tables.py from bench/results/e10_frames.csv
\begin{table}[t]\centering\small
\begin{adjustbox}{max width=\linewidth}\begin{tabular}{lr}
\hline
Quantity & value \\
\hline
frames recorded & 1756 \\
median frame time (ms) & 16.70 \\
frame rate at the median (fps) & 59.9 \\
95th percentile frame time (ms) & 16.80 \\
frames while verifying (iterates arriving) & 1702 \\
\hline
\end{tabular}\end{adjustbox}
\caption{E10: desktop frame times over 30~s of the viewer while a recorded verification stream plays (\texttt{?bench=30}; look thermal, transparency wboit). The browser presents at the display's refresh rate, so the rate is capped by it.}
\label{tab:e10}
\end{table}


\paragraph{Latency of verification (E11).} Table~\ref{tab:e11} times the live server at full
resolution. Both targets, a first band within 100~ms of the prediction and a verified field within
2~s, are missed. The first band waits for the certificate constant $\hat\kappa$; starting its
computation from the previous design's vector, which stays valid because the certificate is
checked again before use, brings the first band down for small edits (the drag row). The verified
time is dominated by about 800 conjugate-gradient iterations on the CPU of the live server; the
batched GPU pipeline of Chapter~\ref{ch:pagp} certifies designs 35 times faster and is the route to
the target, but it is not yet wired into the live server.

% generated by experiments/report_tables.py from bench/results/e11_latency.json, bench/results/e11_latency_drag.json
\begin{table}[t]\centering\small
\begin{adjustbox}{max width=\linewidth}\begin{tabular}{lrrrr}
\hline
designs & $z$ for $\hat\kappa$ & first band (ms) & verified (ms) & commit \\
\hline
10 random & cold & 4956 & 18513 & ea20019 \\
10 drag & warm-started & 856 & 14650 & 76cd736 \\
\hline
\end{tabular}\end{adjustbox}
\caption{E11: median latency of the live server at $256^2\times4$ (FNO ensemble on the GPU, certified PCG on the CPU), from the prediction to the first iterate carrying a guaranteed band and to the verified field ($\tau_v = 0.1$~K). Targets (\S13 of the plan): 100~ms and 2~s; both missed.}
\label{tab:e11}
\end{table}


\section{What was not built}
\begin{itemize}
\item Flux particles and line-integral convolution (T1.2), the blackbody palette with bloom and
tone mapping (T1.3 beyond the lookup tables) and procedural trace generation (T1.4) were cut for
time.
\item In VR, the layer peel and the WebXR session work, but grabbing and moving chips is
desktop-only (the live panel); probe pen, haptics and the copper brush were not built.
\item Uncertainty is shown as a labelled band and a badge, not as the shimmer, hatching or
value-suppressing palettes considered in the plan.
\item The proposer replay (T1.8) was not built, and the measured-overlay feature (T1.9) went with
the hardware validation.
\item The pre-registered pilot study (E12, \texttt{study/prereg.md}) was not run: it needs ethics
approval, which was not obtained in time.
\end{itemize}
